\chapter{76}\label{section-76}

Louis’ restliche Aufstiegszeit verging in null Komma nichts. Sein Gehirn
war so stark gefordert, dass er weder seine müden Beine noch seinen
beschleunigten Puls wahrnahm. Seine Gedanken kreisten pausenlos um
Bastiano. Wenigstens waren sich Valentina und er nicht näher, als Louis
wusste.

Ununterbrochen tauchten die gleichen Fragen auf. Warum war Bastiano im
Visier der Polizei? In welcher Verbindung stand der zu Giorgias Absturz?
Wie waren Spuren von Bastiano an den Absturzort gekommen? Und dann
Valentinas Aussagen zu einem erwiesenen Drogenkonsum bei Giorgia.
Unvorstellbar.

Louis konnte einfach keine Zusammenhänge herstellen. Selbst seine
wildesten Gedankengänge, erdachten Verbindungen und Abläufe, eigentlich
sein ganzes Repertoire an Denkmodellen, versagte.

Und obendrauf wurde er gesucht. Mit Bild. Mit seinen veröffentlichten
Namensinitialen. Es konnte nur eine Frage der Zeit sein, bis einer hier
oben den Aufruf entdeckte. Und als sich Louis auch noch an Valentinas
Worte erinnerte, von wegen \emph{vielleicht stellst du dich am besten,
um dich von deinem Leiden zu befreien} schnürte sich sein Hals vollends
zu. Er musste stehen bleiben. Beugte sich vornüber. Ging schliesslich in
die Knie und starrte den Boden an. Und in diesem Moment wusste er es. Er
würde nichts tun. Gar nichts. Wie all die Tage davor.

Die letzten Meter waren geschafft. Beim Ablegen seines Rucksacks
vibrierte Louis’ Handy. Er setzte sich zerknüllt auf die Bank vor der
Hütte und las eine Nachricht von Joh.

Hi, Ciro. Wann bist du oben? Freu mich auf dich, Joh.

Louis drückte auf das Hörersymbol.

«Joh, hallo. Ich bin gerade angekommen»

«Cool, passt gut. Ich komme rauf, habe in der Hütte etwas zu reparieren.
Mazi und David bleiben bei den Schafen. Kochst du Kaffee?»

«Ja, klar. Mache ich gerne. Bis gleich.»

Louis trug den Rucksack in den kleinen Vorraum der Hütte. Dann suchte er
den Kaffeekocher. Er goss Wasser ein, legte das befüllte Sieb darauf und
schraubte die beiden Teile zusammen. Gedankenverloren schaute er vor
sich hin. Eigentlich war es gerade gut, bald nicht mehr alleine zu sein.

Es dauerte keine Viertelstunde, bis Joh mit erhitztem Gesicht vor ihm
stand. Er strahlte Louis an und nahm ihn in den Arm.

«Schön, dich zu sehen.»

Bei Kaffee und Louis’ mitgebrachten frischen Backwaren löcherte Joh
Louis geradezu mit Fragen zu \emph{Toscas} Zustand. Erzählte, wie es den
Schwarznasen ging und hielt Louis einen Vortrag zur baldigen
\emph{Schafscheid.} Er beschrieb ihm den Abstieg, die Einsätze der
einzelnen Begleiter und \emph{Flexas} Treibaufgabe.

«Wenn wir am Wochenende den ganzen Weg zu Fuss runtergehen, kennst du
die Abstiegsroute, okay?»

Louis nickte.

«Aber jetzt geht’s erstmal um die Dusche,» Joh machte ein ernstes
Gesicht, «die Pumpe klemmt. Ich hol die Werkzeugkiste. Gehst du schon
mal rüber?»
